\section{Описание}

Требуется написать реализацию словаря на основе дерева PATRICIA (Practical Algorithm
To Retrieve Information Coded In Alphanumeric), а затем провести исследование качества
полученной программы.

PATRICIA-дерево~--- это сжатое бинарное trie, в котором каждый узел хранит номер
критического бита ключа, по которому происходит ветвление. Это позволяет избежать
цепочек узлов с одним потомком и хранить только те позиции, в которых ключи
действительно расходятся.

Ключевая особенность структуры~--- \textit{обратные рёбра}: когда при обходе номер
критического бита следующего узла не превышает номера текущего, обход заканчивается.
Узел назначения такого ребра является кандидатом для полного сравнения строк.

Для корректной обработки граничных случаев используется узел-заголовок
(\texttt{header}) с нулевым критическим битом, левый потомок которого является корнем
дерева. Все ключи переводятся в нижний регистр, что обеспечивает
регистронезависимость. Словарь поддерживает операции \texttt{Insert}, \texttt{Find},
\texttt{Remove}, а также бинарные \texttt{Save} и \texttt{Load}.

\pagebreak

\section{Исходный код}

\subsection{Вспомогательные функции}

Функция \texttt{ExtractBit} извлекает $i$-й бит строки (нумерация с 1, старший бит
первого байта~--- бит 1). Функция \texttt{FindFirstDiffBit} находит позицию первого
бита, в котором два ключа различаются,~--- результат используется при вставке для
определения критического бита нового узла. Функция \texttt{NormalizeKey} приводит
ключ к нижнему регистру.

\begin{lstlisting}[language=C++, caption={Вспомогательные функции}]
static int ExtractBit(const std::string &key, int pos) {
    int byteIdx = (pos - 1) / 8;
    int bitOff  = 7 - ((pos - 1) % 8);
    if (byteIdx >= static_cast<int>(key.size())) {
        return 0;
    }
    return (static_cast<unsigned char>(key[byteIdx]) >> bitOff) & 1;
}

static std::string NormalizeKey(const std::string &s) {
    std::string result(s.size(), '\0');
    for (int i = 0; i < static_cast<int>(s.size()); i++) {
        result[i] = static_cast<char>(
            tolower(static_cast<unsigned char>(s[i])));
    }
    return result;
}

static int FindFirstDiffBit(const std::string &a, const std::string &b) {
    int maxLen = static_cast<int>(
        a.size() > b.size() ? a.size() : b.size());
    for (int i = 0; i < maxLen; i++) {
        unsigned char x = (i < static_cast<int>(a.size()))
            ? static_cast<unsigned char>(a[i]) : 0;
        unsigned char y = (i < static_cast<int>(b.size()))
            ? static_cast<unsigned char>(b[i]) : 0;
        if (x != y) {
            unsigned char diff = x ^ y;
            int shift = 0;
            while ((diff & 0x80) == 0) { diff <<= 1; shift++; }
            return i * 8 + shift + 1;
        }
    }
    return maxLen * 8 + 1;
}
\end{lstlisting}

\subsection{Структура узла и класс дерева}

Каждый узел хранит ключ, значение, номер критического бита и указатели на двух
потомков. Поле \texttt{size} в классе \texttt{TPatriciaTrie} отслеживает количество
элементов.

\begin{lstlisting}[language=C++, caption={Объявления TNode и TPatriciaTrie}]
struct TNode {
    std::string key;
    uint64_t value;
    int checkBit;
    TNode *left;
    TNode *right;

    TNode(const std::string &k, uint64_t v, int b);
};

class TPatriciaTrie {
  public:
    TPatriciaTrie();
    ~TPatriciaTrie();

    bool Find(const std::string &key, uint64_t &value) const;
    bool Insert(const std::string &key, uint64_t value);
    bool Remove(const std::string &key);
    bool Save(const std::string &path) const;
    bool Load(const std::string &path);

  private:
    TNode *header;
    int size;

    TNode *SearchNode(const std::string &key) const;
    TNode **FindParentPtr(TNode *target, const std::string &key) const;
    void SaveSubtree(std::ofstream &file, TNode *cur) const;
    void DestroySubtree(TNode *cur);
};
\end{lstlisting}

\subsection{Поиск}

Функция \texttt{SearchNode} спускается по дереву, выбирая ветвь по критическому биту,
до тех пор пока не встретит обратное ребро. Найденный узел является только кандидатом:
\texttt{Find} дополнительно выполняет полное сравнение строк.

\begin{lstlisting}[language=C++, caption={SearchNode и Find}]
TNode *TPatriciaTrie::SearchNode(const std::string &key) const {
    TNode *prev = header;
    TNode *cur = header->left;
    while (cur->checkBit > prev->checkBit) {
        prev = cur;
        if (ExtractBit(key, cur->checkBit)) {
            cur = cur->right;
        } else {
            cur = cur->left;
        }
    }
    return cur;
}

bool TPatriciaTrie::Find(const std::string &key, uint64_t &value) const {
    if (!header) {
        return false;
    }
    TNode *found = SearchNode(NormalizeKey(key));
    if (found->key != NormalizeKey(key)) {
        return false;
    }
    value = found->value;
    return true;
}
\end{lstlisting}

\subsection{Вставка}

При вставке сначала выполняется \texttt{SearchNode}, чтобы найти коллизию или
определить ключ-ориентир. Затем \texttt{FindFirstDiffBit} вычисляет критический бит
нового узла. После этого выполняется повторный спуск до позиции вставки, и новый узел
подвешивается с обратным ребром на самого себя.

\begin{lstlisting}[language=C++, caption={Insert}]
bool TPatriciaTrie::Insert(const std::string &key, uint64_t value) {
    std::string lkey = NormalizeKey(key);
    if (!header) {
        header = new TNode(lkey, value, 0);
        header->left = header;
        header->right = nullptr;
        size++;
        return true;
    }

    TNode *existing = SearchNode(lkey);
    if (existing->key == lkey) {
        return false;
    }

    int diffBit = FindFirstDiffBit(lkey, existing->key);
    TNode *prev = header;
    TNode *cur = header->left;
    while (cur->checkBit > prev->checkBit && cur->checkBit < diffBit) {
        prev = cur;
        if (ExtractBit(lkey, cur->checkBit)) {
            cur = cur->right;
        } else {
            cur = cur->left;
        }
    }

    TNode *newNode = new TNode(lkey, value, diffBit);
    if (ExtractBit(lkey, diffBit)) {
        newNode->right = newNode;
        newNode->left = cur;
    } else {
        newNode->left = newNode;
        newNode->right = cur;
    }

    if (prev == header) {
        header->left = newNode;
    } else if (ExtractBit(lkey, prev->checkBit)) {
        prev->right = newNode;
    } else {
        prev->left = newNode;
    }
    size++;
    return true;
}
\end{lstlisting}

\subsection{Удаление}

Удаление обрабатывает два случая. Если у удаляемого узла есть обратное ребро на самого
себя, родительский указатель переключается на другого потомка. Иначе находится узел,
который содержит обратное ребро на удаляемый элемент; его данные переносятся в
удаляемый узел, а сам найденный узел физически удаляется.

\begin{lstlisting}[language=C++, caption={Remove --- два случая удаления}]
bool TPatriciaTrie::Remove(const std::string &key) {
    if (!header) {
        return false;
    }

    std::string lkey = NormalizeKey(key);
    TNode *target = SearchNode(lkey);
    if (target->key != lkey) {
        return false;
    }

    bool selfRef = (target->left == target || target->right == target);
    if (selfRef) {
        if (target == header) {
            delete header;
            header = nullptr;
            size--;
            return true;
        }
        TNode *other = (target->left == target) ? target->right : target->left;
        *FindParentPtr(target, lkey) = other;
        delete target;
        size--;
        return true;
    }

    /* ... find back-edge holder, move its data, fix links ... */
    size--;
    return true;
}
\end{lstlisting}

\subsection{Сериализация}

Метод \texttt{Save} обходит дерево рекурсивно и записывает каждый узел как тройку:
длина ключа, сам ключ и значение. Длина хранится в \texttt{uint16\_t}, потому что по
условию ключ может иметь длину 256 символов. Метод \texttt{Load} читает записи во
временное дерево, проверяет корректность каждой записи и заменяет текущее состояние
только после успешного чтения всего файла.

\begin{lstlisting}[language=C++, caption={SaveSubtree и Load}]
void TPatriciaTrie::SaveSubtree(std::ofstream &file, TNode *cur) const {
    uint16_t klen = static_cast<uint16_t>(cur->key.size());
    file.write(reinterpret_cast<const char *>(&klen), sizeof(klen));
    file.write(cur->key.data(), klen);
    file.write(reinterpret_cast<const char *>(&cur->value),
               sizeof(cur->value));
    if (cur->left && cur->left->checkBit > cur->checkBit) {
        SaveSubtree(file, cur->left);
    }
    if (cur->right && cur->right->checkBit > cur->checkBit) {
        SaveSubtree(file, cur->right);
    }
}

bool TPatriciaTrie::Load(const std::string &path) {
    std::ifstream file(path, std::ios::binary);
    if (!file) {
        return false;
    }

    TPatriciaTrie tmp;
    while (true) {
        uint16_t klen = 0;
        if (!file.read(reinterpret_cast<char *>(&klen), sizeof(klen))) {
            if (!file.eof()) { return false; }
            break;
        }
        if (klen == 0 || klen > MAX_KEY_LENGTH) {
            return false;
        }
        std::string key(klen, '\0');
        if (!file.read(&key[0], klen)) { return false; }

        uint64_t val = 0;
        if (!file.read(reinterpret_cast<char *>(&val), sizeof(val))) {
            return false;
        }
        if (!tmp.Insert(key, val)) {
            return false;
        }
    }
    std::swap(header, tmp.header);
    std::swap(size, tmp.size);
    return true;
}
\end{lstlisting}

\subsection{Обработка команд}

Функция \texttt{main} читает команды построчно. Добавление, удаление, поиск,
сохранение и загрузка выводят ответы в формате условия.

\begin{lstlisting}[language=C++, caption={Обработка команд в main}]
while (std::getline(std::cin, line)) {
    std::istringstream ss(line);
    std::string cmd;
    ss >> cmd;
    if (cmd == "+") {
        std::string word;
        uint64_t val;
        ss >> word >> val;
        std::cout << (trie.Insert(word, val) ? "OK" : "Exist") << "\n";
    } else if (cmd == "-") {
        std::string word;
        ss >> word;
        std::cout << (trie.Remove(word) ? "OK" : "NoSuchWord") << "\n";
    } else if (cmd == "!") {
        std::string subcmd, path;
        ss >> subcmd >> path;
        if (subcmd == "Save") {
            std::cout << (trie.Save(path) ? "OK" : "ERROR: cannot save") << "\n";
        } else if (subcmd == "Load") {
            std::cout << (trie.Load(path) ? "OK" : "ERROR: cannot load") << "\n";
        }
    } else {
        uint64_t val;
        if (trie.Find(cmd, val)) {
            std::cout << "OK: " << val << "\n";
        } else {
            std::cout << "NoSuchWord\n";
        }
    }
}
\end{lstlisting}

\pagebreak

\section{Консоль}

\begin{alltt}
root@container:/workspace# cmake -B build -DCMAKE_BUILD_TYPE=Release
root@container:/workspace# cmake --build build
root@container:/workspace# ./build/lab2/lab2_main
+ a 1
OK
+ A 2
Exist
A
OK: 1
- A
OK
a
NoSuchWord
\end{alltt}

\pagebreak
